\newenvironment{bprooftree}
{\leavevmode\hbox\bgroup}
{\DisplayProof\egroup}

We highlight some features of the dependent type theory that we work in, which is given in its entirety in \autoref{appendix:type-theory}. 
Our starting point is Martin-L\"of type theory with the empty type (\(\bot\)), unit type (\(\top\)) and boolean type (\(\sf{Bool}\)), together with dependent pair types (\(\Sigma\)), dependent function types (\(\Pi\)), and propositional equality types (\(=\)).
For \(\Sigma\) and \(\Pi\) we include definitional \(\eta\) rules.

We also include a Tarski-style universe \(U\), meaning that it comes with an explicit decoder \( \El \) from terms of \(U\) to types, which we omit when possible for readability.
The universe contains terms that act as codes for the basic types, and similar codes for \(\Pi x:A.\,B(x)\), \(\Sigma x:A.\,B(x)\) and \(a=_Aa'\) when both \(A\) and \(B(x)\) have codes in \(U\).
For each type (former), we use a superscript \((-)^U\) to distinguish the code (term of \(U\)) from the type when needed, although we will often omit this for readability.
We will call types with codes in the universe \emph{small}, whereas general types are \emph{large}.
A consequence of the $\eta$-rule for $\Pi$-types is that we may identify a family of small types $x:A\vdash B(x):\U$ indexed by the type $A$ with a function $B:A\to\U$; we will make extensive use of this in later sections.

Although we do not work with univalence, we will make use of terminology and typical constructions from homotopy type theory \cite{hottbook}.
In particular, the notion of an equivalence will be important to this work: we would like to say that a function \(f:A\to B\) is an \emph{equivalence}, written \(f:A\simeq B\), if it has a \emph{quasi-inverse}, which is a map \(g:B\to A\) such that both compositions are pointwise equal to the identity. However, we want to ensure that the type stating that \(f\) is an equivalence a \emph{mere proposition}, meaning that all terms of the type are equal. To this end, we also assume a coherence condition between the proofs of pointwise equality; see \cite[Definition 4.2.1]{hottbook}, where this definition is called a half-adjoint equivalence. We say that a type $A$ is \emph{contractible} when there exists exactly one term of type $A$ up to propositional equality.
We may also refer to terms of identity types as paths, and use operations defined by path induction, such as the action of a function \( f : A \to B \) on paths and the transport operation.

We also include the \emph{function extensionality} axiom as described in homotopy type theory \cite{hottbook}, which states that for any two dependent functions \(f,g:\Pi x:A.\,B(x)\), the canonical function \(\mathsf{happly} : (f=g)\to\Pi x:A.\,(f\,x=g\,x)\) defined by path induction is an equivalence: \[
    (f=g)\simeq \Pi{x:A} \. (f\,x=g\,x).\tag{function extensionality}
\]
We call the quasi-inverse to \(\mathsf{happly}\) obtained via this equivalence \(\mathsf{funext}\).

\subsection{Sizes and parametric quantifiers}
\label{ss:sizes}
We extend the theory with a primitive type \( \Size \), which comes with a zero term and a successor operation: \begin{center}
    \begin{bprooftree}
        \AxiomC{\(\phantom\Gamma\)}
        %\RightLabel{\(\sf{Size}\)-form,}
        \UnaryInfC{$\Gamma\vdash\Size\Type$}
    \end{bprooftree} \quad
    \begin{bprooftree}
        \AxiomC{\(\phantom\Gamma\)}
        %\RightLabel{\(\sf{Size}\)-intro\(_0\),}
        \UnaryInfC{$\Gamma\vdash 0:\Size$}
    \end{bprooftree} \quad
    \begin{bprooftree}
        \AxiomC{$\Gamma\vdash i:\Size$}
        %\RightLabel{\(\sf{Size}\)-intro\(_\sf{suc}\).}
        \UnaryInfC{$\Gamma\vdash{\uparrow}i:\Size$}
    \end{bprooftree}
\end{center}
So far this looks a lot like the type of natural numbers; the difference in their behaviour will be the elimination principle, for which we first need to introduce some other notions.
In particular, we assume a preorder (also known as quasi-order) relation \(\leq\) on the type of sizes, which will be a small type.
We encode this relation as a mere proposition, which means that for any two terms \( p, q : i \leq j \) we have a propositional equality \( p = q \).
The ordering on sizes is encoded by the following rules:
\begin{center}
    \begin{bprooftree}
        \AxiomC{$\Gamma\vdash i:\Size$}
        \AxiomC{$\Gamma\vdash j:\Size$}
        %\RightLabel{\(\le\)-form,}
        \BinaryInfC{$\Gamma\vdash i\le j\Type$}
    \end{bprooftree}
    \begin{bprooftree}
        \AxiomC{$\Gamma \vdash i : \Size$}
        %\RightLabel{\(\le\)-intro\(_0\),}
        \UnaryInfC{$\Gamma \vdash{\sf{le}_0}\,i:0 \le i$}
    \end{bprooftree}\quad
    \begin{bprooftree}
        \AxiomC{$\Gamma \vdash i : \Size$}
        %\RightLabel{\(\le\)-intro\(_{\sf{suc}}\),}
        \UnaryInfC{$\Gamma \vdash{\sf{le}_{\sf{suc}}}\,i:i \le {\uparrow}i $}
    \end{bprooftree} \\[2ex]
    \begin{bprooftree}
        \AxiomC{$\Gamma \vdash i : \Size$}
        %\RightLabel{\(\le\)-intro\(_{\sf{refl}}\),}
        \UnaryInfC{$\Gamma \vdash\sf{le}_{\sf{refl}}\,i:i \le i$}
    \end{bprooftree}\quad
    \begin{bprooftree}
        \AxiomC{$\Gamma \vdash p:i\le j$}
        \AxiomC{$\Gamma \vdash q:j\le k$}
        %\RightLabel{\(\le\)-intro\(_{\sf{trans}}\).}
        \BinaryInfC{$\Gamma \vdash\sf{le}_{\sf{trans}}\,p\,q:i \le k$}
    \end{bprooftree}
\end{center}
We let $i < j$ be notation for $\uparrow \!i \le j$. 
To discuss the elimination and computation principles for sizes, we first need to add universal and existential quantifiers over sizes.
These are impredicative: although the type of sizes is large, quantifying over small types will result in a small type:
\begin{center}
    \begin{bprooftree}
        \AxiomC{$\Gamma, i : \Size \vdash A(i) : U$}
        \UnaryInfC{$\Gamma \vdash \exists i.\,A(i) : U$}
    \end{bprooftree}\quad\begin{bprooftree}
        \AxiomC{$\Gamma, i : \Size \vdash A(i) : U$}
        \UnaryInfC{$\Gamma \vdash \forall i.\,A(i) : U$}
    \end{bprooftree}
\end{center}
These quantifiers will have very similar rules to \(\Sigma\) and \(\Pi\); however, we will introduce additional axioms in \autoref{ss:equivalences} to internalise a notion of parametricity: sizes must be treated uniformly by quantifiers.
As with \(\Sigma i:\Size.\,A(i)\), terms of \(\exists i.\,A(i)\) are pairs \( \langle s, a \rangle ^\exists\) where \( s: \Size \) and \( a : A(s) \).
However, unlike \(\Sigma\), the elimination principle of \(\exists\) only allows elimination to a small type \(P\):
\begin{center}
    \begin{bprooftree}
        \AxiomC{$\Gamma\vdash s:\Size$}
        \AxiomC{$\Gamma\vdash a : A$}
        \BinaryInfC{$\Gamma\vdash\<s,a\>^\exists:\exists i. A(i)$}
    \end{bprooftree} \quad
    \begin{bprooftree}
        \AxiomC{$\Gamma, z : \exists i. A(i) \vdash P(z) : \U$}
        \noLine
        \UnaryInfC{$\Gamma, i : \Size, x : A \vdash p(i,x) :P(\langle i, x \rangle^\exists)$}
        \UnaryInfC{$\Gamma, z : \exists i. A(i) \vdash \ind_\exists(z, p(i,x)) :P(z)$}
    \end{bprooftree} \\[2ex]
    \begin{bprooftree}
        \AxiomC{}
        \UnaryInfC{$\Gamma\vdash\ind_\exists(\<s,a\>^\exists, p(i,x))\equiv p(s,a) : {P(\<s,a\>^\exists)}$}
    \end{bprooftree}
\end{center}
Consequently, we cannot define a projection to the first component, since the type \( \Size \) is not small.
Thus, the size given to the second component of a term of the existential type is inaccessible via the elimination principle and therefore kept abstract.
%This notion of the existential quantifier is closer to the classical notion, where existence does not necessarily give access to a witness.
For \(\forall\), the rules are completely analogous to \(\Pi\):
\begin{center}
    \begin{bprooftree}
        \AxiomC{$\Gamma, i : \Size \vdash a(i) : A(i)$}
        \UnaryInfC{$\Gamma \vdash \lambda^\forall i : \Size \. a(i): \forall i . A(i) $}
    \end{bprooftree}\quad
    \begin{bprooftree}
        \AxiomC{$\Gamma \vdash f : \forall i. A(i)$}
        \AxiomC{$\Gamma \vdash s : \Size$}
        \BinaryInfC{$\Gamma \vdash f \, s : A(s)$}
    \end{bprooftree}
\end{center}
\begin{center}
    \begin{bprooftree}
        \AxiomC{}
        \UnaryInfC{$\Gamma \vdash (\lambda^\forall i \: \Size \. a(i)) \, s \equiv a(s):A(s)$}
    \end{bprooftree}\quad
    \begin{bprooftree}
        \AxiomC{}
        \UnaryInfC{$\Gamma \vdash (\lambda^\forall i \: \Size \. f \, i) \equiv f :\forall i . A(i)$}
    \end{bprooftree}
\end{center}
Since we will not allow case distinctions on sizes due to its limited elimination principle, it will not be possible to construct non-parametric functions in the syntax.

Here we have used \(\<\dots\>^\exists\) and \(\lambda^\forall\) to distinguish the pairs and functions from those in \(\Sigma\)-types and \(\Pi\)-types; however, when this does not lead to ambiguity we will drop these superscripts.
In fact, we see \(\Pi i:\Size.\,A(i) \simeq \forall i.\,A(i) \) where the maps are given by \(h\mapsto\lambda^\forall i.\,h\,i\) and \(h'\mapsto\lambda^\Pi i.\,h'\,i\).
In particular, this shows that function extensionality for \(\Pi\) implies a function extensionality principle for \(\forall\).
However, the dual principle \(\exists i.\,A(i)\simeq\Sigma i.\,A(i)\) can not be shown, since the elimination principle of \(\exists\) does not allow us to build a function to \(\Sigma\), and this equivalence is not true in the model of \autoref{sec:model}.

% Note that function extensionality for \(\Pi\) implies a function extensionality principle for \(\forall\): for \(f,g:\forall i.\,A(i)\) we can construct a quasi-inverse of the canonical map \((f=g)\to\forall i.\,(f\,i=g\,i)\); to do this, assume that we we have \(\forall i.\,(f\,i=g\,i)\), then with function extensionality for \(\Pi\) we see that \(\lambda^\Pi i.\,f\,i=\lambda^\Pi i.\,g\,i\), and therefore by applying the map \((\Pi i:\Size.\,A(i))\to(\forall i.\,A(i))\) given by \(h\mapsto\lambda^\forall i.\,h\,i\) we get \(
%     f\equiv\lambda^\forall i.\,(\lambda^\Pi j.\,f\,j)\,i=\lambda^\forall i.\,(\lambda^\Pi j.\,g\,j)\,i\equiv g.
% \)

Another important notion will be `bounded' quantification over sizes.
Recall that for sizes \( i, j : \Size \), the type \( i < j \) is a mere proposition. We will use the notation $i : \Size \vdash \forall j < i. A(j)$ to denote the type $i : \Size \vdash \forall j. (j < i) \to A(j)$. Similarly $i : \Size \vdash \exists j < i. A(i)$ will denote the type $i : \Size \vdash \exists j. ((j < i) \times A(j))$. When defining functions in the type theory, we will often omit the inequality proofs, writing the term $\langle j, \langle p, a \rangle \rangle : \exists j < i. A(j)$ as $\langle j, a\rangle $ instead, as \( p \) is unique up to propositional equality. Similarly, if we have a term $f : \forall j < i \. A(j)$, we will simply write $f \, j$ to mean $f \, j \, p$ where $p: j < i$. In lambda notation, we may also write $\lambda j < i$ when constructing a term of type $\forall j < i\. A(j)$. 

Finally, we can give the elimination principle for the type $\Size$, which comes in the form of a fixpoint operator, corresponding to well-founded induction on sizes:
    \begin{center}
        \begin{bprooftree}
            \AxiomC{$\Gamma, i : \Size \vdash A(i) \Type$}
            \AxiomC{$\Gamma \vdash f : \forall i . ((\forall j < i. A(j)) \to A(i))$}
            \BinaryInfC{$\Gamma \vdash \fix f : \forall i. A(i)$}
        \end{bprooftree}
    \end{center}
The $\beta$-rule for $\fix$ ensures that \(\sf{fix}\,f\) behaves as a fixpoint of the map \(g\mapsto \lambda i.\,f\,i\,(\lambda j<i.\,g\,i)\):
    \begin{center}
    \begin{bprooftree}
            \AxiomC{$\Gamma \vdash f : \forall i . ((\forall j < i. A(j)) \to A(i))$}
            \UnaryInfC{$\Gamma \vdash \fix^\beta \, f : \forall i \. \fix f \, i = f \, i \, (\lambda j < i\. \fix f  \, j)$}
        \end{bprooftree}
    \end{center}
We can view this as the unfolding rule of the fixpoint operator.
The model we construct supports this \(\beta\)-rule as a definitional equality, because the model is extensional, meaning that propositional and definitional equality coincide.
However, this definitional equality would lead to a loss of decidable type checking if allowed to unfold indefinitely.
Thus, we opt to encode this as a propositional equality for now, and discuss alternatives in \autoref{sec:conclusion}.
Importantly, we can prove that, up to propositional equality, $\fix \, f$ is the \emph{unique} term satisfying this unfolding, which we may view as its propositional $\eta$-rule:
\begin{lemma}[Uniqueness of fixpoints]\label{lemma:unique-fix}
The following type is contractible for any term $f :  \forall i.(\forall j< i. A(j)) \to A(i)$:
\[
\sum_{h :\forall i .A(i)}  \forall i \. (h \, i = f \, i \, (\lambda j < i\. h \, j)).
\]
\end{lemma}
\begin{proof}
We follow a well-known strategy \cite[Theorem 5.4.7]{hottbook} \cite[Lemma 3.2]{guarded-cubical}.
First, note that the type is inhabited by the pair $(\fix f, \fix^\beta \, f)$. For any other pair $(g, g^\beta)$, we need to show $(\fix f, \fix^\beta f) = (g, g^\beta)$.
For this, it suffices to give a pointwise equality \(e:\forall i.\fix f\,i=g\,i\) and a proof \(e^\beta:(\sf{funext}(e))_*\,(\fix^\beta f)=g^\beta\).
We take \(e\coloneqq \fix(\lambda i\.\lambda e'\.p)\) where \(p\) is a concatenation of \[\hfill\begin{tikzcd}
    \fix f\,i \ar[r,"\fix^\beta f\,i",equal] & f\,i\,(\lambda j<i\.\fix f\,j) \ar[rrr,"{\mathsf{ap}_{fi}(\mathsf{funext}(\lambda j<i\.e'\,j))}",equal] &&& f\,i\,(\lambda j<i\.g\,j) \ar[r, "g^\beta\,i",equal] & g\,i.
\end{tikzcd}\hfill\]
To define \(e^\beta\), it suffices to show that the following square commutes: \[\hfill\begin{tikzcd}
    \fix f\,i \ar[r,"\fix^\beta f\,i",equal]\ar[d,"e\,i"',equal] & f\,i\,(\lambda j<i\.\fix f\,j) \ar[d,"{\mathsf{ap}_{fi}(\mathsf{funext}(\lambda j<i\.e\,j))}",equal] \\
    g\,i & f\,i\,(\lambda j<i\.g\,j). \ar[l, "g^\beta\,i",equal]
\end{tikzcd}\hfill\]
However, this is given precisely by the \(\beta\)-rule \(\fix^\beta\,(\lambda i\.\lambda e'\.p)\).
\end{proof}

\subsection{Type equivalences}
\label{ss:equivalences}
We want to internalise parametricity for quantifying over sizes in our theory, which we do by adding additional equivalences to the theory as axioms. These equivalences describe how universal and existential quantification over sizes commute with the type constructors for small types. They will allow us to encode (co)inductive types in the calculus, and their consistency will be justified by the model we construct in \ref{sec:model}. 

Before introducing the additional axioms for the calculus, we show that two equivalences can be constructed in the calculus already.
These correspond to swapping function arguments and elements of pairs. 
\begin{proposition}\label{prop:definable-eqvs}
For small types \(A\) and \(x:A,i:\sf{Size}\vdash B(x,i)\) we have the following equivalences:
\begin{align*}
    \Pi x : A. \, \forall i. B(x, i)&\simeq\forall i \. \Pi x : A.\,B(i, x), \tag{\(\exists\) parametric over \(\Sigma\)} \\
    \Sigma x : A. \, \exists i . B(x, i)&\simeq\exists i \. \Sigma x : A. \, B(i, x). \tag{\(\forall\) parametric over \(\Pi\)}
\end{align*}
\end{proposition}
We also require equivalences showing the remaining interactions between the quantifiers and \(\Sigma\), \(\Pi\), and \(<\), for all \emph{small types} \(A\), \(x:A,i:\sf{Size}\vdash B(x,i)\), and \(i: \Size \vdash C(i)\).
These are not constructible in the calculus, so we add these as axioms, similarly to what has been done in guarded type theory \cite{mogelberg, guardedtypes16}.
\begin{align*}
    %A &\simeq \forall i.\, A &
    %\exists i \. A &\simeq A &
    %i \notin \text{fv}(A) \\
    \Pi x : A.\,\exists i . B(x, i)&\simeq\exists i \. \Pi x : A. \, \exists j < i .B(x, j), \tag{\(\exists\) parametric over \(\Pi\)} \\
    \Sigma x:A.\,\forall i . B(x, i) &\simeq \forall i \. \Sigma x : A.\,B(i, x), \tag{\(\forall\) parametric over \(\Sigma\)} \\
    \exists i \.  C(i)&\simeq \exists i \. \exists j < i \.C(j), \tag{\(\exists\) parametric over \(<\)} \\
    \forall i \.  C(i) &\simeq \forall i \. \forall j < i \.C(j).  \tag{\(\forall\) parametric over \(<\)}
\end{align*}
More precisely, for each equivalence, a \emph{canonical} map $f$ can already be defined in one direction: for \(\exists\) these are the right-to-left maps and for \(\forall\) these are the left-to-right maps. 
Our axioms state that each such $f$ forms an equivalence, giving us a quasi-inverse in the other direction.

Note that the first equivalence differs in the fact that the existential quantifier does not exactly commute with the small $\Pi$-type; we will call this \emph{weakly commuting} later on. The intuitive reason for this is as follows: to go from right to left, there must be a limit of all sizes in the range of the dependent product. However, we cannot guarantee that the type $B$ is monotone (also known as covariant) in $\Size$, so we make use of $\exists j < i .B(x, j)$ instead, which is monotone in $\Size$. Importantly, this axiom thus provides an implicit construction of limits of sequences of sizes in the theory, without explicitly adding a limit operator on sizes to the syntax.

With these parametricity axioms, we can also show how the quantifiers interact with fixed types such as the base types: 
\begin{proposition}
    For a small type \(A\) with \(i\notin\sf{fv}(A)\), we have the following: \begin{align*}
        \exists i.\,A\simeq A, \tag{\(\exists\) parametric over a small fixed type} \\
        \forall i.\,A\simeq A. \tag{\(\forall\) parametric over a small fixed type}
    \end{align*}
\end{proposition}
\begin{proof}
    We first prove both statements in the special case where \(A\coloneqq\top\).
    For \(\exists\), we use the fact that for any type \(B\) we have \(\bot\to B\simeq\top\), and for \(\forall\), we use function extensionality: \begin{align*}
        \exists i.\,\top
        \simeq\exists i.(\bot\to\exists j<i.\,\bot)
        \simeq\bot\to\exists i.\,\bot
        \simeq \top, &&
        \forall i.\,\top\simeq\top.
    \end{align*}
    We use this special case and the equivalence \(\Sigma x:A.\,\top\simeq A\) to prove the general versions: \begin{gather*}
        \exists i.\,A\simeq\exists i.\,\Sigma x:A.\,\top\simeq\Sigma x:A.\,\exists i.\,\top\simeq\Sigma x:A.\,\top\simeq A, \\
        \forall i.\,A\simeq\forall i.\,\Sigma x:A.\,\top\simeq\Sigma x:A.\,\forall i.\,\top\simeq\Sigma x:A.\,\top\simeq A.\qedhere
    \end{gather*}
\end{proof}
% \begin{figure}[ht]
%     For small types \(A\), \(B(i,x)\), and \(C(i)\) with \(i \notin \text{fv}(A)\), 
%     \begin{align*}
%         %A &\simeq \forall i.\, A &
%         %\exists i \. A &\simeq A &
%         %i \notin \text{fv}(A) \\
%         \Sigma x:A.\,(\forall i . B(x, i)) &\simeq \forall i \. \Sigma x : A.\,B(i, x), &
%         \exists i \. \Pi x : A. \, \exists j < i .B(x, j) &\simeq \Pi x : A.\,\exists i . B(x, i), \\        
%         \forall i \.  C(i) &\simeq \forall i \. \forall j < i \.C(j), &
%         \exists i \. \exists j < i \.C(j) &\simeq \exists i \.  C(i).
%     \end{align*}
%     \caption{Type equivalences for existential and universal quantification. For each equivalence, a left-to-right map is definable. The additional axiom states that this map forms an equivalence.} 
%     \label{type-equivs}
% \end{figure}
% % Context Formers
The rules for context formation are given as follows:
    \begin{center}
        \begin{bprooftree}
            \AxiomC{$\phantom{\Gamma, x : A \vdash M : B}$}
            \UnaryInfC{$\vdash \diamond \; \Ctxt$}
        \end{bprooftree}\quad
        \begin{bprooftree}
            \AxiomC{$\vdash \Gamma \; \Ctxt$}
            \AxiomC{$\Gamma \vdash A \Type$}
            \BinaryInfC{$\vdash \Gamma, x : A \; \Ctxt$}
        \end{bprooftree}
    \end{center}
    \begin{center}
        \begin{bprooftree}
            \AxiomC{$\vdash \Gamma, x : A, \Delta \; \Ctxt$}
            \UnaryInfC{$\Gamma, x:A , \Delta \vdash x:A$}
        \end{bprooftree}\quad
        \begin{bprooftree}
            \AxiomC{$\Gamma, x : A\vdash b(x):B(x)$}
            \AxiomC{$\Gamma\vdash a:A$}
            \BinaryInfC{$\Gamma\vdash b(a):B(a)$}
        \end{bprooftree}
    \end{center}
\vspace{10pt}

The rules for the empty-, unit-, and boolean type are as follows, where their elimination and computation principles are left implicit.
% empty, unit, bool
\begin{center}
    \begin{bprooftree}
        \AxiomC{$\phantom{\Gamma \vdash M : \Pi x : A.\,B}$}
        \UnaryInfC{$\Gamma \vdash \bot \Type$}
    \end{bprooftree}\quad
    \begin{bprooftree}
        \AxiomC{$\phantom{\Gamma \vdash M : \Pi x : A.\,B}$}
        \UnaryInfC{$\Gamma \vdash \top \Type$}
    \end{bprooftree}\quad
    \begin{bprooftree}
        \AxiomC{$\phantom{\Gamma \vdash M : \Pi x : A.\,B}$}
        \UnaryInfC{$\Gamma \vdash \mathsf{Bool} \Type$}
    \end{bprooftree}\quad
\end{center}
\begin{center}
    \begin{bprooftree}
        \AxiomC{\phantom{$\Gamma, x : A \vdash b : B$}}
        \UnaryInfC{$\Gamma \vdash \star : \top$}
    \end{bprooftree}\quad
    \begin{bprooftree}
        \AxiomC{\phantom{$\Gamma, x : A \vdash b : B$}}
        \UnaryInfC{$\Gamma \vdash \mathsf{tt} : \mathsf{Bool}$}
    \end{bprooftree}\quad
    \begin{bprooftree}
        \AxiomC{\phantom{$\Gamma, x : A \vdash b : B$}}
        \UnaryInfC{$\Gamma \vdash \mathsf{ff} : \mathsf{Bool}$}
    \end{bprooftree}\quad
\end{center}
\vspace{10pt}

The rules for $\Pi$-types include definitional $\beta$- and $\eta$-rules.
% Pi type
\begin{center}
    \begin{bprooftree}
        \AxiomC{$\Gamma \vdash A \Type$}
        \AxiomC{$\Gamma, x : A \vdash B(x) \Type$}
        \BinaryInfC{$\Gamma \vdash \Pi x : A.\,B(x) \Type$}
    \end{bprooftree}\quad
\end{center}
\begin{center}
    \begin{bprooftree}
        \AxiomC{$\Gamma, x : A \vdash b(x) : B(x)$}
        \UnaryInfC{$\Gamma \vdash \lambda x \: A \. b(x): \Pi x : A.\,B(x)$}
    \end{bprooftree}\quad
    \begin{bprooftree}
        \AxiomC{$\Gamma \vdash f : \Pi x : A.\,B(x)$}
        \AxiomC{$\Gamma \vdash a : A$}
        \BinaryInfC{$\Gamma \vdash f \,a : B(a)$}
    \end{bprooftree}
\end{center}
\begin{center}
    \begin{bprooftree}
        \AxiomC{$\Gamma, x : A \vdash b(x) : B(x)$}
        \AxiomC{$\Gamma \vdash a : A$}
        \BinaryInfC{$\Gamma \vdash (\lambda x \: A \. b(x)) \, a \equiv b(a): B(a)$}
    \end{bprooftree}\quad
    \begin{bprooftree}
        \AxiomC{$\Gamma \vdash f : \Pi x : A.\,B(x)$}
        \UnaryInfC{$\Gamma \vdash (\lambda x \: A \. f \, x) \equiv f : \Pi x : A.\,B(x)$}
    \end{bprooftree}
\end{center}
\vspace{10pt}

Similarly to $\Pi$-types, $\Sigma$-types also come with $\beta$- and $\eta$-rules. The $\eta$-rule for $\Sigma$-types is defined via its destructors, where $\pi_1$ and $\pi_2$ are the projections derived from the elimination principle. 
\begin{center}
    \begin{bprooftree}
        \AxiomC{$\Gamma \vdash A \Type$}
        \AxiomC{$\Gamma, x : A \vdash B(x) \Type$}
        \BinaryInfC{$\Gamma \vdash \Sigma x:A.\,B(x) \Type$}
    \end{bprooftree}\quad
    \begin{bprooftree}
        \AxiomC{$\Gamma \vdash a : A$}
        \AxiomC{$\Gamma \vdash b : B(a)$}
        \BinaryInfC{$\Gamma \vdash \< a, b \> : \Sigma x:A.\,B(x)$}
    \end{bprooftree}
    \end{center}
    \begin{center}
        \begin{bprooftree}
            \AxiomC{$\Gamma, z : \Sigma x:A.\,B(x) \vdash P(z)\Type$}
            \noLine
            \UnaryInfC{$\Gamma, x : A, y : B \vdash p(x,y) : P(\< x, y \>)$}
            \UnaryInfC{$\Gamma, z : \Sigma x:A.\,B(x) \vdash \ind_\Sigma(z, p(x,y)) : P(z)$}
        \end{bprooftree}
    \end{center}
    \begin{center}
        \begin{bprooftree}
            \AxiomC{$\Gamma, z : \Sigma x:A.\,B(x) \vdash P(z)\Type$}
            \noLine
            \UnaryInfC{$\Gamma, x : A, y : B \vdash p : P(\< x, y \>)$}
            \UnaryInfC{$\Gamma, x : A, y : B \vdash \ind_\Sigma(\< x, y \>, p) \equiv p(x,y) : P(\< x, y \>)$}
        \end{bprooftree}
    \end{center}
    \begin{center}
        \begin{bprooftree}
            \AxiomC{$\Gamma \vdash c : \Sigma x:A.\,B(x)$}
            \UnaryInfC{$\Gamma \vdash c \equiv \< \pi_1(c), \pi_2(c) \> : \Sigma x:A.\,B(x)$}
        \end{bprooftree}
    \end{center}
\vspace{10pt}

Propositional identity types are defined with the standard induction principle.
% Identity Type
    \begin{center}
        \begin{bprooftree}
            \AxiomC{$\Gamma \vdash A \Type$}
            \AxiomC{$\Gamma \vdash a_1, a_2 : A$}
            \BinaryInfC{$\Gamma \vdash a_1 =_A a_2 \Type$}
        \end{bprooftree}\quad
        \begin{bprooftree}
            \AxiomC{$\Gamma \vdash A \Type$}
            \AxiomC{$\Gamma \vdash a : A$}
            \BinaryInfC{$\Gamma \vdash \refl \, a : a =_A a$}
        \end{bprooftree}\quad
    \end{center}
    \begin{center}
        \begin{bprooftree}
            \AxiomC{$\Gamma , x : A, y: A, z : x =_A y \vdash P(x, y, z) \Type$}
            \noLine
            \UnaryInfC{$\Gamma , x : A \vdash p(x) : P(x, x, \refl \, x)$}
            \UnaryInfC{$\Gamma, x:A, y:A, z: x=_A y \vdash \ind_=(x,y, z, p) : P(x, y, z)$}
        \end{bprooftree}
    \end{center}
    \begin{center}
        \begin{bprooftree}
            \AxiomC{$\Gamma , x : A, y : A, z : x =_A y \vdash P(x, y, z) \Type$}
            \noLine
            \UnaryInfC{$\Gamma , x : A \vdash p(x) : P(x, x, \refl \, x)$}
            \UnaryInfC{$\Gamma, x:A \vdash \ind_=(x, x, \refl \, x, p) \equiv p(x) : P(x, x, \refl \, x)$}
        \end{bprooftree}
    \end{center}
\vspace{10pt}

The universe type is closed under the base types, propositional equality, and the type constructors. The definitional type equalities for the small types ensure that they lift to the proper type.
    \begin{center}
        \begin{bprooftree}
        \AxiomC{\phantom{$\Gamma \vdash u : \U$}}
        \UnaryInfC{$\Gamma \vdash \U \Type$}
        \end{bprooftree}\quad
        \begin{bprooftree}
        \AxiomC{$\Gamma \vdash A : \U$}
        \UnaryInfC{$\Gamma \vdash \El A \, \Type$}
        \end{bprooftree}
    \end{center}
    \begin{center}
        \begin{bprooftree}
            \AxiomC{$\phantom{\Gamma \vdash}$}
            \UnaryInfC{$\Gamma \vdash \bot^\U : \U$}
        \end{bprooftree}\quad
        \begin{bprooftree}
            \AxiomC{$\phantom{\Gamma \vdash}$}
            \UnaryInfC{$\Gamma \vdash \top^\U : \U$}
        \end{bprooftree}\quad
        \begin{bprooftree}
            \AxiomC{$\phantom{\Gamma \vdash}$}
            \UnaryInfC{$\Gamma \vdash \mathsf{Bool}^\U : \U$}
        \end{bprooftree}
    \end{center}
    \begin{center}
        \begin{bprooftree}
            \AxiomC{$\Gamma \vdash A : U$}
            \AxiomC{$\Gamma, x : A \vdash B(x) : \U$}
            \BinaryInfC{$\Gamma \vdash \Pi^U x:A.\, B(x): \U$}
        \end{bprooftree}
    \end{center}
    \begin{center}
        \begin{bprooftree}
            \AxiomC{$\Gamma \vdash A : \U$}
            \AxiomC{$\Gamma, x : \El A \vdash B(x) : \U$}
            \BinaryInfC{$\Gamma \vdash \Sigma^\U x:A.\, B(x) : \U$}
        \end{bprooftree}
    \end{center}
    \begin{center}
        \begin{bprooftree}
            \AxiomC{$\Gamma \vdash A : \U$}
            \AxiomC{$\Gamma\vdash a : \El A$}
            \AxiomC{$\Gamma\vdash a' : \El A$}
            \TrinaryInfC{$\Gamma \vdash a=^U_Aa' : \U$}
        \end{bprooftree}
    \end{center}
%Type eqs
    \begin{center}
    \begin{align*}
        \El (\bot^\U) &\equiv \bot & \El (\Sigma^U x:A.\, B(x)) &\equiv \Sigma_{x:\El A} (\El B(x)) \\
        \El (\top^\U) &\equiv \top & \El (\Pi^U x:A.\, B(x)) &\equiv \Pi_{x:\El A} (\El B(x)) \\
        \El (\mathsf{Bool}^\U) &\equiv \mathsf{Bool} & \El(a=^U_Aa')&\equiv(a=_{\El A}a')
    \end{align*}
    \end{center}
\vspace{10pt}

The type of sizes is defined with constructors for zero and successor. The ordering on sizes is defined as a mere proposition, with constructors for reflexivity and transitivity.
% Sizes
\begin{center}
    \begin{bprooftree}
        \AxiomC{\phantom{$\Gamma \vdash s_1, s_2 : \Size$}}
        \UnaryInfC{$\Gamma \vdash \Size \Type$}
    \end{bprooftree}\quad
    \begin{bprooftree}
        \AxiomC{$\Gamma \vdash  i, j : \Size$}
        \UnaryInfC{$\Gamma \vdash \; i \leq j : U$}
    \end{bprooftree}\quad
    \begin{bprooftree}
        \AxiomC{\phantom{$\Gamma \vdash s_1, s_2 : \Size$}}
        \UnaryInfC{$\Gamma \vdash 0 : \Size$}
    \end{bprooftree}\quad
    \begin{bprooftree}
        \AxiomC{$\Gamma \vdash i  : \Size$}
        \UnaryInfC{$\Gamma \vdash {\uparrow \! i} : \Size$}
    \end{bprooftree}
\end{center}

\begin{center}
    \begin{bprooftree}
        \AxiomC{$\Gamma\vdash i:\Size$}
        \AxiomC{$\Gamma\vdash j:\Size$}
        %\RightLabel{\(\le\)-form,}
        \BinaryInfC{$\Gamma\vdash i\le j\Type$}
    \end{bprooftree}
    \begin{bprooftree}
        \AxiomC{$\Gamma \vdash i : \Size$}
        %\RightLabel{\(\le\)-intro\(_0\),}
        \UnaryInfC{$\Gamma \vdash{\sf{le}_0}\,i:\El(0 \le i)$}
    \end{bprooftree}\quad
    \begin{bprooftree}
        \AxiomC{$\Gamma \vdash i : \Size$}
        %\RightLabel{\(\le\)-intro\(_{\sf{suc}}\),}
        \UnaryInfC{$\Gamma \vdash{\sf{le}_{\sf{suc}}}\,i:\El(i \le {\uparrow}i) $}
    \end{bprooftree} \\[2ex]
    \begin{bprooftree}
        \AxiomC{$\Gamma \vdash i : \Size$}
        %\RightLabel{\(\le\)-intro\(_{\sf{refl}}\),}
        \UnaryInfC{$\Gamma \vdash\sf{le}_{\sf{refl}}\,i:\El(i \le i)$}
    \end{bprooftree}\quad
    \begin{bprooftree}
        \AxiomC{$\Gamma \vdash p:i\le j$}
        \AxiomC{$\Gamma \vdash q:j\le k$}
        %\RightLabel{\(\le\)-intro\(_{\sf{trans}}\).}
        \BinaryInfC{$\Gamma \vdash\sf{le}_{\sf{trans}}\,p\,q:\El(i \le k)$}
    \end{bprooftree} \\[2ex]
    \begin{bprooftree}
        \AxiomC{$\Gamma \vdash p:i\le j$}
        \AxiomC{$\Gamma \vdash q:i\le j$}
        \BinaryInfC{$\Gamma \vdash\sf{le}_{\sf{eq}}\,p\,q :\El(p=q)$}
    \end{bprooftree}
\end{center}
\vspace{10pt}

The type of the existential quantifier over sizes is defined as follows. Note that the elimination principle is limited to small types.
    \begin{center}
        \begin{bprooftree}
            \AxiomC{$\Gamma, i : \Size \vdash A(i) : \U$}
            \UnaryInfC{$\Gamma \vdash \exists i . A(i) : \U $}
        \end{bprooftree}
    \end{center}
    \begin{center}
        \begin{bprooftree}
            \AxiomC{$\Gamma, i : \Size \vdash A(i) : \U$}
            \AxiomC{$\Gamma \vdash s : \Size$}
            \AxiomC{$\Gamma \vdash a : \El A(i)$}
            \TrinaryInfC{$\Gamma \vdash \langle s, a \rangle^\exists : \El (\exists i . A(i))$}
        \end{bprooftree}
    \end{center}
    \begin{center}
        \begin{bprooftree}
            \AxiomC{$\Gamma, z : \El (\exists i. A(i)) \vdash P(z) : \U$}
            \noLine
            \UnaryInfC{$\Gamma, i : \Size, x : \El A \vdash p(i,x) : \El P(\langle i, x \rangle^\exists)$}
            \UnaryInfC{$\Gamma, z : \El(\exists i.A(i))\vdash \ind_\exists(z,p(i,x)) : \El P(z)$}
        \end{bprooftree}
    \end{center}
    \begin{center}
        \begin{bprooftree}
            \AxiomC{$\Gamma, z : \El (\exists i. A(i)) \vdash P(z) : \U$}
            \noLine
            \UnaryInfC{$\Gamma, i : \Size, x : \El A \vdash p : \El P(\langle i, x \rangle^\exists)$}
            \noLine
            \UnaryInfC{$\Gamma\vdash s:\Size$\qquad$\Gamma\vdash a:\El A(s)$}
            \UnaryInfC{$\Gamma\vdash \ind_\exists(\langle s, a \rangle^\exists, p(i,x)) \equiv p(s,a) : \El P(\langle s,a \rangle^\exists)$}
        \end{bprooftree}
    \end{center}
\vspace{10pt}

The universal quantifier is defined analogously to the $\Pi$-type, only as an impredicative quantifier.
\begin{center}
    \begin{bprooftree}
        \AxiomC{$\Gamma, i : \Size \vdash A(i) :U$}
        \UnaryInfC{$\Gamma \vdash \forall i . A(i) :U$}
    \end{bprooftree}\quad
\end{center}
\begin{center}
    \begin{bprooftree}
        \AxiomC{$\Gamma, i : \Size \vdash a(i) : A(i)$}
        \UnaryInfC{$\Gamma \vdash \lambda^\forall i : \Size \. a(i): \El (\forall i . A(i)) $}
    \end{bprooftree}\quad
    \begin{bprooftree}
        \AxiomC{$\Gamma \vdash f : \El(\forall i . A(i))$}
        \AxiomC{$\Gamma \vdash s : \Size$}
        \BinaryInfC{$\Gamma \vdash f \, s : \El A(s)$}
    \end{bprooftree}
\end{center}
\begin{center}
    \begin{bprooftree}
        \AxiomC{$\Gamma, i : \Size \vdash a(i) : \El A(i)$}
        \AxiomC{$\Gamma \vdash s : \Size$}
        \BinaryInfC{$\Gamma \vdash (\lambda^\forall i \: \Size \. a(i)) \, s \equiv a(s): \El A(s)$}
    \end{bprooftree}\quad
    \begin{bprooftree}
        \AxiomC{$\Gamma \vdash f : \El(\forall i . A(i))$}
        \UnaryInfC{$\Gamma \vdash (\lambda^\forall i \: \Size \. f \, i) \equiv f : \El (\forall i . A(i))$}
    \end{bprooftree}
\end{center}
\vspace{10pt}

Finally, the fixpoint operator on sizes, corresponding to well-founded induction, is defined as follows. Note that the associated $\beta$-rule is propositional.
    \begin{center}
        \begin{bprooftree}
            \AxiomC{$\Gamma, i : \Size \vdash A(i) \Type$}
            \AxiomC{$\Gamma \vdash f : \forall i . ((\forall j < i. A(j)) \to A(i))$}
            \BinaryInfC{$\Gamma \vdash \fix f : \forall i. A(i)$}
        \end{bprooftree}
    \end{center}
    \begin{center}
    \begin{bprooftree}
            \AxiomC{$f : \forall i . ((\forall j < i. A(j)) \to A(i))$}
            \UnaryInfC{$\Gamma \vdash \fix^\beta \, f : \forall i \. \fix f \, i = f \, i \, (\lambda j < i\. \fix f  \, j)$}
        \end{bprooftree}
    \end{center}
